\section{Introduction}

Memes are among the most prolific forms of digital communication, yet they remain deeply challenging
for multimodal AI. Unlike image-caption pairs where visual and textual content reinforce each other,
memes deliberately exploit \emph{incongruity}—a serene image paired with a desperate caption, a
cheerful face accompanying nihilistic text. This cross-modal conflict is not noise to be filtered; it is
the semantic mechanism through which humor, sarcasm, and irony are constructed. As a result, standard
multimodal models that assume or enforce cross-modal agreement systematically fail on meme affect.

The challenge compounds when we ask: \emph{how should a meme sound?} Human narrators instinctively
adjust pitch, pace, and emotional coloring when reading meme text aloud—leaning into irony,
exaggerating deadpan delivery, or escalating expressive intensity to match the visual absurdity. Current
text-to-speech (TTS) systems have no access to this pragmatic context; they generate audio from text
semantics alone, missing the affective layer encoded in the image-text relationship. This produces audio
that is semantically correct but expressively hollow—natural in isolation, inappropriate for the meme.

A third, underexplored dimension is the audio modality itself: once a voiced meme exists, does the
prosodic rendering add information beyond what image and text already communicate? Do different
emotional renderings of the same script change how viewers perceive sentiment and intention? These
questions have practical implications for accessible content, meme retrieval, and affective computing,
yet no prior work has systematically studied tri-modal (image, text, audio) interaction in the meme domain.

\textbf{In this paper}, we introduce \textbf{MemeSonic}, a unified framework that addresses meme
understanding and expressive audio generation jointly. Our contributions are:

\begin{enumerate}
  \item \textbf{Explicit Affect Decomposition.} We decompose meme understanding into independent
    per-modality emotion distributions and compute a continuous \emph{incongruity score} that
    quantifies cross-modal conflict. This turns the incongruity—previously treated as noise—into an
    explicit, differentiable signal available throughout the pipeline.

  \item \textbf{Incongruity-Aware MMOE Fusion.} We propose a Mixture-of-Experts fusion mechanism
    whose router is conditioned on the incongruity signal. Three expert strategies handle distinct
    interaction types: Redundancy (Concat), Uniqueness (Late Fusion), and Synergy (Incongruity
    Fusion). This is fundamentally different from standard MoE systems that route on task difficulty
    or data statistics; our router encodes the \emph{semantic relationship} between modalities.

  \item \textbf{Meme-to-Speech Pipeline.} We translate the
    fused affective representation into expressive audio via an adapter that conditions TTS synthesis
    on the incongruity score, enabling prosodic variance to reflect not just the text sentiment but
    the cross-modal pragmatic intent of the meme.

  \item \textbf{Tri-Modal Alignment Study.} We conduct a controlled experiment across a 3$\times$3
    prosody grid (stability $\times$ style) to quantify whether audio modality contributes independent
    affective information beyond image-text fusion. Using Emotion2Vec embeddings and a learned linear
    projector to bridge modality spaces, we provide the first systematic evidence that (a) audio
    carries non-redundant emotional signal (KL divergence: $V$\,$\leftrightarrow$\,$A$ = 1.14 vs.\
    $V$\,$\leftrightarrow$\,$T$ = 0.69), and (b) the failure of naive fusion is a \emph{modality gap
    problem}, not an absence of information.
\end{enumerate}

Across 10 classification tasks on three benchmarks (MET-Meme, Memotion 7k, MMSD 2.0), our
Incongruity-Aware Fusion achieves the best Macro-F1 on tasks dominated by non-literal language
(metaphor +87.6\%, humour +12.1\% over Late Fusion). Human evaluation of generated audio
(n=84 ratings) shows MemeSonic achieves AMOS 3.95 vs.\ 1.68 for Vanilla TTS, with 84.5\%
preference over the vanilla baseline—demonstrating that incongruity-aware conditioning substantially
improves pragmatic appropriateness without sacrificing naturalness.
